% 2015 年重庆卷（理工农医类）第 7 题：程序框图（依固定扫描页重绘，只画拓扑与标号，不按比例）
\begin{tikzpicture}[>=stealth,scale=0.85]
  % 开始
  \draw[rounded corners=6pt] (-1.0,11.6) rectangle (1.0,12.4);
  \node at (0,12.0) {开始};
  \draw[->] (0,11.6)--(0,11.2);
  % 赋初值
  \draw (-1.8,10.4) rectangle (1.8,11.2);
  \node at (0,10.8) {$s=0,\ k=0$};
  \draw[->] (0,10.4)--(0,8.1);
  % 判断框（空白）
  \draw (-1.9,7.4)--(0,6.7)--(1.9,7.4)--(0,8.1)--cycle;
  % 否
  \draw[->] (0,6.7)--(0,6.0);
  \node at (0.35,6.35) {否};
  \draw (-1.4,6.0)--(1.4,6.0)--(1.2,5.2)--(-1.6,5.2)--cycle;
  \node at (-0.1,5.6) {输出 $k$};
  \draw[->] (0,5.2)--(0,4.6);
  \draw[rounded corners=6pt] (-1.0,4.6) rectangle (1.0,3.8);
  \node at (0,4.2) {结束};
  % 是：先 k=k+2，再 s=s+1/k，回到判断框上方
  \draw[->] (1.9,7.4)--(4.6,7.4)--(4.6,7.9);
  \node at (2.4,7.78) {是};
  \draw (3.5,7.9) rectangle (5.7,8.7);
  \node at (4.6,8.3) {$k=k+2$};
  \draw[->] (4.6,8.7)--(4.6,9.1);
  \draw (3.4,9.1) rectangle (5.8,9.9);
  \node at (4.6,9.5) {$s=s+\frac{1}{k}$};
  \draw[->] (4.6,9.9)--(4.6,10.2)--(0,10.2);
\end{tikzpicture}
